
\documentclass[12pt,a4paper]{article}
\usepackage[utf8]{inputenc}
\usepackage[english]{babel}
\usepackage{enumerate}
\usepackage{hyperref}

\title{Master's Thesis: Combined Research Proposals and Integrated Introduction}
\author{Filip Rosiak \& Jakub Stefański}
\date{February 2026}

\begin{document}

\maketitle

\section{Proposed Research Topics (General)}
These topics aim to bridge the gap between discrete structural analysis and continuous latent optimization for complex 3D constructions.

\begin{enumerate}
    \item \textbf{Synergistic Linkage Learning:} Combining Discrete Epistasis Estimation and Continuous Latent Mapping for 3D Construction Optimization.
    \item \textbf{Learning and Exploiting Structural Hierarchies:} Advanced Evolutionary Synthesis of Complex 3D Modular Figures.
    \item \textbf{From Discrete Substructures to Continuous Latent Spaces:} A Multi-Representation Framework for 3D Construction Design.
    \item \textbf{Context-Aware Optimization of 3D Modular Robots:} Evaluating Implicit and Explicit Dependency Learning.
    \item \textbf{Navigating the Morphological Landscape:} Integrated Discrete Measures and Latent Search for 3D Construction Synthesis.
    \item \textbf{Beyond Padding:} Robust Evolutionary Optimization of Variable-Length 3D Constructions via Neural Representations.
    \item \textbf{Hybrid Model-Based Evolutionary Design:} Coupling CoSo and RV-GOMEA for Efficient 3D Figure Synthesis.
    \item \textbf{Assessing Structural Neutrality:} A Comparative Study of Geometric Metrics and Neural Embeddings in 3D Design.
    \item \textbf{Automated Identification of Functional Modules:} Enhancing 3D Construction Evolution through Linkage Discovery.
    \item \textbf{Neural-Evolutionary Architectures:} Optimized Synthesis of 3D Modular Constructions in Learned Spaces.
    \item \textbf{Quantifying and Reducing Epistasis in Morphological Search:} Techniques for Complex 3D Figure Design.
    \item \textbf{Multi-Scale Dependency Mapping:} Learning the Grammar and Geometry of 3D Modular Constructions.
    \item \textbf{Evolutionary Synthesis of 3D Figures:} Harnessing Latent Manifolds and Contextual Substructure Swaps.
    \item \textbf{Structural Integrity and Locomotion:} Co-optimizing Form and Function in 3D Modular Robotics.
    \item \textbf{Transfer Learning in 3D Design:} Utilizing Learned Embeddings to Accelerate Evolutionary Optimization.
    \item \textbf{Interpretable Latent Spaces for 3D Construction Design:} Mapping Geometric Measures to Continuous Dimensions.
    \item \textbf{Scalable Evolutionary Synthesis:} Multi-Level Linkage Learning for Large-Scale 3D Modular Figures.
    \item \textbf{Generative Models for 3D Constructions:} Bridging Symbolic Genotypes and Geometric Latent Spaces.
    \item \textbf{Epistasis-Aware Genetic Programming:} Enhancing 3D Construction Design with Real-Valued GOMEA and Neutral Measures.
    \item \textbf{Autonomous Morphogenesis:} A Combined Framework for Learning Structural Dependencies in 3D Modular Figures.
\end{enumerate}

\section{Integrated Introduction}

The automated synthesis of complex 3D modular constructions represents a frontier in evolutionary computation and robotics. These constructions, often composed of interconnected modules with varying morphologies, are typically represented by variable-length, non-positional genotypes. A fundamental challenge in this domain is the presence of high epistasis—where the functional contribution of a genotype fragment depends heavily on its context. Standard evolutionary operators frequently fail to account for these dependencies, resulting in the destructive disruption of successful modular building blocks.

To overcome these challenges, this research proposes a unified framework that integrates two complementary approaches to structural optimization. The first approach focuses on the \textbf{explicit estimation of epistasis} within the discrete domain. By developing and evaluating local geometric, structural, and textual measures, we aim to identify "neutral" substructure swaps—swaps that preserve the functional integrity of the construction. This allows algorithms like Context-aware Substructure Optimization (CoSo) to perform more efficient searches by respecting the natural dependencies within the 3D construction's grammar.

The second approach utilizes \textbf{deep generative models} to map these complex, discrete genotypes into a continuous, stałowymiarową (fixed-dimensional) latent space. By training Variational Autoencoders (VAEs) and their grammar-aware variants, we can transform the variable-length search space into a dense, continuous manifold. This enables the use of Real-Valued Gene-pool Optimal Mixing Evolutionary Algorithms (RV-GOMEA), which can implicitly learn the relationships between morphological parameters through advanced linkage learning mechanisms.

The synergy between these two methods allows for a comprehensive understanding of the 3D construction's design space. Explicit discrete measures provide local, interpretable insights into structural neutrality, while latent space mappings offer global, continuous tools for exploring complex nieliniowe (non-linear) dependencies. Together, they form a robust toolkit for the automated design of high-performance 3D modular constructions, capable of evolving both form and behavior in a highly efficient manner.

\section{Experimental Framework}
The proposed combined research will involve several key experimental phases:
\begin{enumerate}
    \item \textbf{Benchmark Comparison:} Evaluating the efficiency of CoSo-guided discrete searches versus RV-GOMEA-guided latent searches across diverse 3D construction objectives (e.g., maximization of velocity, structural stability, or biomass efficiency).
    \item \textbf{Cross-Domain Correlation Analysis:} Investigating the relationship between discrete geometric neutrality measures and the proximity of solutions within the learned latent space.
    \item \textbf{Representation Learning Sensitivity:} Analyzing how the choice of autoencoder architecture (e.g., standard VAE vs. Grammar-VAE) affects the quality and "evolvability" of the latent manifold for 3D figures.
    \item \textbf{Hybrid Operator Evaluation:} Testing a hybrid approach where discrete measures identify candidate modules that are then refined using continuous optimization in the latent space.
    \item \textbf{Epistasis Reduction Study:} Quantifying the degree of epistasis reduction achieved by both explicit measures and latent mappings, and its impact on the scalability of the evolutionary process.
\end{enumerate}

\end{document}
